\chapter{The render context and the pipeline}
\label{ch:pipeline}

\begin{aside}
We have talked about \code{view} returning an element, layout
assigning rects, and the renderer emitting bytes. The piece in
between --- the \emph{pipeline} that ties them together --- is
worth a chapter of its own. It is where \maya{} decides what to
recompute and what to reuse, and it is the boundary between the
pure world of \code{update}/\code{view} and the messy world of
file descriptors.
\end{aside}

\section{What the pipeline does}

Each frame, the pipeline:

\begin{enumerate}[leftmargin=*]
  \item Asks the program for a fresh \code{Element} via
    \code{view(model)}.
  \item Lays it out against the current terminal size.
  \item Paints it into a canvas of packed cells.
  \item Diffs the canvas against the previous one.
  \item Serialises the diff into ANSI bytes.
  \item Writes the bytes to the terminal.
\end{enumerate}

Six steps. Each step has a well-defined input and output. The
boundaries are stable; the implementations can change. That is
the discipline of a pipeline architecture.

\subsection*{Why a pipeline?}

Two alternatives existed:

\begin{itemize}[leftmargin=*]
  \item \textbf{Direct rendering.} \code{view} would return a
    sequence of ``draw'' calls that the terminal executes
    directly. Tight coupling. Tested only at the integration
    level. Imagine debugging \emph{any} part of it.
  \item \textbf{Single-pass rendering.} Merge layout and paint
    into one walk. Faster in principle. But you lose the ability
    to inspect intermediate results, and adding features
    (themes, hit testing, animations) requires touching the one
    monolithic pass.
\end{itemize}

A pipeline gives you intermediate values you can inspect, test,
and replace. Layout takes an element and produces a layout tree.
Paint takes a layout tree and produces a canvas. Diff takes two
canvases and produces a region list. Each is a function from one
value to another; each is testable in isolation.

\begin{idea}
A pipeline is what you get when you apply the Elm architecture's
``effects as data'' philosophy to the renderer itself. The
intermediate values are the data; the stages are the functions.
\end{idea}

\section{The RenderContext}

\file{maya/include/maya/core/render\_context.hpp} defines the
struct that flows through the pipeline:

\begin{lstlisting}
struct RenderContext {
    Element       tree;        // from view()
    LayoutTree    layout;      // after measure/distribute/place
    Canvas        canvas;      // after paint
    DiffResult    diff;        // diff against previous canvas
    SerializeResult bytes;     // ANSI bytes to write
    Stats         stats;       // timings, counts
};
\end{lstlisting}

Each field is populated by one stage. By the end of the frame,
all are filled. The context is then reused for the next frame
(only \code{canvas} and \code{stats} persist; the rest are
recomputed).

\subsection*{Why a single struct?}

We could have used six free functions chained by hand. Bundling
into a struct does three things:

\begin{itemize}[leftmargin=*]
  \item \textbf{Logging.} The context can be dumped at any
    point. ``Why is this frame slow?'' is answerable by
    inspecting \code{stats}; ``why is this cell wrong?'' by
    inspecting \code{canvas}.
  \item \textbf{Lifetimes.} The struct owns all its fields. Pass
    a \code{const RenderContext\&} into any stage and it can
    read whatever it needs without lifetime puzzles.
  \item \textbf{Threading.} If we ever want to run layout on a
    different thread from paint, the context is the unit of
    handoff.
\end{itemize}

\section{The six stages, expanded}

\subsection*{Stage 1: view}

The pipeline calls \code{P::view(model)}. The result is an
\code{Element}. This step is allocator-pressure heavy but
otherwise trivial: it is your code.

Cost: depends on \code{view}. Microseconds for a typical UI.

\subsection*{Stage 2: layout}

Convert the \code{Element} tree to a \code{LayoutTree}: a flat
\code{vector<LayoutNode>} where each node has a final
\code{Rect}. Layout is its own three-phase pass (measure,
distribute, place) detailed in Chapter~\ref{ch:layout}.

Cost: O(n) in tree size. Tens of microseconds for hundreds of
nodes.

\subsection*{Stage 3: paint}

Walk the layout tree and write packed cells into a \code{Canvas}.
A \code{PackedCell} is a 64-bit struct holding a codepoint, a
style ID (an index into a pool), and flags. The canvas is a flat
\code{vector<PackedCell>} of width $\times$ height entries.

Cost: O(cells) = O(width $\times$ height). For 80$\times$24,
that is 1920 entries. Microseconds.

\subsection*{Stage 4: diff}

Compare the current \code{Canvas} against the previous one. The
output is a list of \code{Row} differences: for each row that
changed, the leftmost and rightmost dirty cell.

The diff is SIMD-accelerated (Chapter~\ref{ch:simd}). For an
80$\times$24 grid, it is sub-microsecond on modern CPUs.

\subsection*{Stage 5: serialize}

Convert the diff into ANSI bytes. This stage minimises:

\begin{itemize}[leftmargin=*]
  \item Cursor moves: only emit when the cursor would not
    naturally arrive at the next dirty cell.
  \item SGR sequences: only emit when the style changes from the
    previous cell.
  \item Buffer reuse: the byte buffer carries over capacity from
    frame to frame.
\end{itemize}

Cost: O(dirty cells). For a typical frame with light changes, a
few thousand bytes.

\subsection*{Stage 6: write}

Write the bytes to the terminal via the writer abstraction
(\file{maya/include/maya/terminal/writer.hpp}). The writer
batches calls and flushes once per frame. Synchronous; the cost
is the kernel write syscall plus terminal latency.

Cost: dominated by kernel time. Tens of microseconds typically.

\section{Caching and the cache\_id}

Some pipeline stages can be skipped if their inputs have not
changed. \maya{} tracks input identity with \code{cache\_id} ---
a 64-bit fingerprint computed from the input.

\begin{lstlisting}
struct CacheId {
    uint64_t value;
    bool operator==(const CacheId&) const = default;
};
\end{lstlisting}

If the new \code{CacheId} for the element tree matches the
previous frame's, the entire pipeline can be skipped: no work,
no bytes. The frame is silent.

This is the optimisation that lets a 60 Hz idle app burn
essentially no CPU. \code{view} re-runs every frame (cheap), but
if it produces the same tree, no paint, no diff, no write.

\subsection*{Computing the fingerprint}

The fingerprint walks the tree once, mixing every leaf's content
into a hash. Two trees with identical content produce identical
fingerprints. Two trees that differ in even one cell produce
different fingerprints (with overwhelming probability; collisions
are extremely rare).

\maya{} uses xxHash3 (fast, non-cryptographic) for the mixing.
The walk is itself O(n) in tree size --- same as layout --- so the
trade-off is one extra walk to skip up to five.

\begin{aside}
This is the same trick React uses with \code{shouldComponentUpdate}
or the newer \code{React.memo}: skip work if the inputs are
unchanged. The difference: \maya{}'s cache key is computed
automatically from the tree, not declared by the user.
\end{aside}

\subsection*{Damage tracking}

Some changes do not affect the visual output. A model field that
\code{view} does not consult can change without producing a
different element. Cache fingerprints catch this for free: the
tree is the same, so the fingerprint is the same.

Other changes affect only a sub-region. If a single line changes
in a status bar, only that row's diff is non-empty. The
serializer emits a cursor move plus the changed cells; everything
else is silent.

\section{Stats and observability}

The \code{Stats} struct holds per-frame timings and counters:

\begin{lstlisting}
struct Stats {
    Duration view_time;
    Duration layout_time;
    Duration paint_time;
    Duration diff_time;
    Duration serialize_time;
    Duration write_time;
    Duration total;

    size_t   cells_total;
    size_t   cells_dirty;
    size_t   bytes_written;
    bool     full_redraw;
};
\end{lstlisting}

After each frame the pipeline emits the stats; if you wire a
debugging widget to read them, you can see exactly where time
goes. \file{maya/include/maya/widget/} ships a
\code{render\_stats} widget for this purpose.

\begin{idea}
The pipeline is observable because the context is a value.
Logging, tracing, profiling, replaying --- all become easy when
the data flows through a single typed channel.
\end{idea}

\section{Adding a stage}

Suppose you want to add a stage between paint and diff that
applies a CRT scanline effect (every other row darkened). The
pipeline accommodates this:

\begin{enumerate}[leftmargin=*]
  \item Add a new field to \code{RenderContext}:
    \code{Canvas filtered\_canvas;}
  \item Insert the stage:
    \code{filtered\_canvas = scanlines(canvas);}
  \item Change the diff input from \code{canvas} to
    \code{filtered\_canvas}.
\end{enumerate}

No other stage cares. The pipeline structure makes the change
local.

\maya{} does not ship a scanline filter, but the example shows
how the pipeline can be extended. Real extensions in flight:
gamma correction for image widgets, optional anti-aliased text.

\section{Threading model}

\maya{}'s pipeline is single-threaded on the main loop. The
reasons:

\begin{itemize}[leftmargin=*]
  \item \textbf{Determinism.} A single thread guarantees
    deterministic ordering, which simplifies replay and tests.
  \item \textbf{Latency.} The pipeline is fast enough that
    threading would add more overhead than it saves. A 60 Hz
    UI has 16ms per frame; \maya{}'s pipeline takes under 1ms.
  \item \textbf{Simplicity.} No locks, no thread-safety
    requirements on \code{view}, no race conditions.
\end{itemize}

\code{Cmd::task} and \code{Cmd::isolated\_task} use worker
threads, but they return their results as messages onto the main
loop's queue. The pipeline itself never runs on a worker thread.

\begin{aside}
If you ever need to render off-thread (e.g. a background
preview), you can run the pipeline as a free function on any
thread you like --- the stages are pure. \code{maya::run}'s
pipeline is single-threaded; the algorithm is not.
\end{aside}

\section{Pitfalls}

\begin{warning}
\textbf{Doing work in view that should be done in update.} If
your \code{view} performs an expensive computation (parsing JSON,
walking a graph), it pays that cost \emph{every frame}. Cache the
result in \code{Model} and update it only when the source data
changes.
\end{warning}

\begin{warning}
\textbf{Large element trees.} The cache fingerprint is fast but
not free. A tree with tens of thousands of nodes will spend
non-trivial time on the fingerprint walk. Use \code{when} and
\code{visible} to elide branches you do not need.
\end{warning}

\begin{warning}
\textbf{Misreading stats.} Frame total is dominated by
\code{write\_time} on most terminals; that is the terminal's
latency, not yours. Focus on \code{view\_time} +
\code{layout\_time} for things you can optimise.
\end{warning}

\begin{warning}
\textbf{Forcing redraw too often.} \code{Cmd::force\_redraw}
discards the cache and emits every cell. Useful after a
disturbance, abusive every frame. Use it only when you must.
\end{warning}

\section{What this means in practice}

You will rarely write code that touches the pipeline directly.
\code{maya::run} sets it up; you write \code{init}, \code{update},
\code{view}, \code{subscribe} and the pipeline handles the rest.

But knowing the pipeline exists lets you:

\begin{itemize}[leftmargin=*]
  \item Reason about performance. ``My frame is slow'' becomes
    ``which stage is slow?''.
  \item Reason about cost. ``Why is my CPU at 5\% on an idle
    app?'' is answered by ``the cache is not catching the
    no-change case''.
  \item Trust that the pipeline is testable. Each stage is a
    pure function from one value to another.
\end{itemize}

\section*{Workshop}

\begin{enumerate}[leftmargin=*]
  \item Find \file{maya/include/maya/render/pipeline.hpp} and
    identify the six stages. Note that the source closely
    mirrors the description in this chapter.
  \item Run a \maya{} app with the debug stats widget enabled
    (set \code{MAYA\_RENDER\_STATS=1}). Watch \code{view\_time},
    \code{layout\_time}, and \code{bytes\_written} change as
    you interact.
  \item Hypothesise: which stage will dominate if you render a
    full-screen text editor with 10\,000 lines? Try it and
    confirm.
\end{enumerate}

\begin{recap}
The pipeline is the six-stage transformation from \code{Model} to
ANSI bytes: view, layout, paint, diff, serialize, write. The
\code{RenderContext} struct carries the intermediates. Cache
fingerprints skip the entire pipeline on no-change. The result is
predictable latency, observable performance, and a clean place to
add new stages.
\end{recap}
